% ─────────────────────────────────────────────────────────────────────────────
%  Capítulo 8 – Asignación de Registros (Chaitin-Briggs)
% ─────────────────────────────────────────────────────────────────────────────
\chapter{Asignación de registros: Chaitin-Briggs}
\label{chap:registro}

\section{Motivación}

Los registros de la CPU son escasos y de acceso mucho más rápido que la memoria.
La \emph{asignación de registros} consiste en asignar las variables y temporales
de la \IR{} a registros físicos (o a posiciones en la pila si no hay
registros suficientes: \emph{spilling}).

El algoritmo de \textbf{Chaitin-Briggs}~\cite{chaitin1982} modela el
problema como un \emph{problema de coloración de grafos}: dos variables
que están \emph{vivas} en el mismo punto del programa no pueden compartir
el mismo registro, y por tanto se conectan en el \emph{grafo de interferencias}.
Colorear el grafo con $K$ colores equivale a asignar $K$ registros.

\section{Registros disponibles}

Se reservan los siguientes registros \emph{caller-saved} para temporales,
evitando tener que salvarlos en llamadas a función:

\begin{table}[H]
  \centering
  \caption{Registros asignables por el compilador}
  \label{tab:regs}
  \begin{tabular}{ll}
    \toprule
    \textbf{Clase} & \textbf{Registros} ($K = 6$) \\
    \midrule
    Enteros  & \reg{rcx}, \reg{rdx}, \reg{r8}, \reg{r9}, \reg{r10}, \reg{r11} \\
    Flotantes & \reg{xmm2}, \reg{xmm3}, \reg{xmm4}, \reg{xmm5}, \reg{xmm6}, \reg{xmm7} \\
    \bottomrule
  \end{tabular}
\end{table}

\section{Algoritmo}

El algoritmo consta de tres fases:

\subsection{Fase 1: Análisis de vivacidad}

Se calculan los conjuntos \emph{use} y \emph{def} por bloque básico y
se resuelven las ecuaciones de punto fijo:
\[
  \text{live\_in}[B]  = \text{use}[B]  \cup (\text{live\_out}[B] \setminus \text{def}[B])
\]
\[
  \text{live\_out}[B] = \bigcup_{S \in \text{succ}(B)} \text{live\_in}[S]
\]

\begin{lstlisting}[style=pseudo, caption={Análisis de vivacidad (punto fijo)}]
Entrada: bloques B_1,...,B_n con conjuntos use[B], def[B]
Salida:  live_in[B], live_out[B] para cada bloque

repetir hasta que no haya cambios:
    para cada bloque B (en orden inverso):
        live_out[B] = union{ live_in[S] | S sucesor de B }
        live_in[B]  = use[B] union (live_out[B] menos def[B])
\end{lstlisting}

\subsection{Fase 2: Grafo de interferencias}

Se recorre cada bloque básico de forma inversa manteniendo el conjunto
$\text{live}$ de variables vivas en el punto actual. Al definir una
variable $d$ se añade una arista $(d, v)$ para cada $v \in \text{live}$.

\begin{lstlisting}[style=cpp, caption={Construcción del grafo de interferencias}]
for (auto it = block.rbegin(); it != block.rend(); ++it) {
    const auto& q = *it;
    if (q.result.kind == IRArg::TEMP || q.result.kind == IRArg::VAR) {
        int d = q.result.addr;
        for (int live_var : live)
            graph.addEdge(d, live_var);  // d interfiere con live_var
        live.erase(d);                   // d deja de estar vivo antes de def
    }
    // Add uses to live set
    if (q.arg1 es variable/temporal) live.insert(q.arg1.addr);
    if (q.arg2 es variable/temporal) live.insert(q.arg2.addr);
}
\end{lstlisting}

\subsection{Fase 3: Coloración del grafo}

Se sigue el algoritmo \emph{simplify}:

\begin{enumerate}
  \item Mientras exista un nodo con grado $< K$: apilarlo y eliminarlo del grafo.
  \item Si todos los nodos tienen grado $\geq K$: elegir el de mayor grado para
    \emph{spill} y eliminarlo (se marca como derramado).
  \item Repetir hasta vaciar el grafo.
  \item Desapilar asignando colores (registros) no usados por vecinos ya coloreados.
\end{enumerate}

\begin{lstlisting}[style=cpp, caption={Coloración simplificada de Chaitin-Briggs}]
// Fase simplify: apila nodos con grado < K
stack<int> nodeStack;
while (!worklist.empty()) {
    auto it = find_if(worklist, [&](int n){ return degree[n] < K; });
    if (it != worklist.end()) {
        nodeStack.push(*it);
        removeNode(*it);
    } else {
        // spill: eliminar nodo de mayor grado
        int spill = *max_element(worklist, [&](int a, int b){
            return degree[a] < degree[b]; });
        spilledNodes.insert(spill);
        removeNode(spill);
    }
}
// Fase assign: asignar colores
while (!nodeStack.empty()) {
    int n = nodeStack.top(); nodeStack.pop();
    set<string> usedColors;
    for (int neighbor : adj[n])
        if (color.count(neighbor)) usedColors.insert(color[neighbor]);
    for (const auto& reg : AVAILABLE_REGS)
        if (!usedColors.count(reg)) { color[n] = reg; break; }
}
\end{lstlisting}

\section{Uso del resultado}

El mapa de coloración \texttt{RegMap} (dirección IR → nombre de registro)
se pasa al \texttt{X86Generator}. Las variables con \texttt{color = ""}
(derramadas) se mantienen en la pila y se acceden con la notación
\texttt{-N(\%rbp)}.

\begin{verbatim}
$ ./compiler --reg-alloc tests/backend/t04_while.java 2>&1
=== Asignacion de Registros ===
  addr 1 (i) -> %rcx
  addr 2 (s) -> %rdx
  addr 3 (t0) -> %r8
  addr 4 (t1) -> spill
\end{verbatim}
